\begin{frame}{실무 규칙: early stopping이 필수}
  \begin{columns}[T]
    \begin{column}{0.46\textwidth}
      \begin{itemize}
        \item \textbf{early stopping이 필수}
      \end{itemize}
      \vskip0.6em
      \begin{block}{표준 워크플로}
        \small \textbf{``learning\_rate를 작게 하고, n\_estimators를
        early stopping으로 결정''} — 실무 표준.
      \end{block}
    \end{column}
    \begin{column}{0.5\textwidth}
      \centering
      \includegraphics[width=\linewidth]{ch07_3_rf_vs_gbdt_curve.png}
      \vskip0.3em
      {\footnotesize RF는 트리를 늘려도 검증 성능이 평탄, GBDT는 정점을
      지나면 다시 나빠진다.}
    \end{column}
  \end{columns}


\note{\begin{itemize}
\item \texttt{n\_estimators}는 RF처럼 ``많을수록(안정적으로) 좋다''가
      \textbf{아님} — 검증 곡선의 \textbf{정점(최저 MSE)을 찾아 거기서
      멈추는 early stopping이 필수}(sklearn \texttt{validation\_fraction}
      옵션으로 내장)
\item \texttt{learning\_rate}를 \textbf{작게}(0.1 $\to$ 0.01) 하면 각
      스텝이 작아져 같은 수준까지 \textbf{더 많은 단계}가 필요 $\to$
      정점의 위치·높이가 \textbf{완만해져 early stopping이 정밀}해짐
\item 이 ``정점을 지나 다시 올라가는 곡선'' = Week 1.1 \textbf{편향-분산
      트레이드오프}: 트리 적으면(불충분 학습) \textbf{편향}이 크고,
      많으면(과적합) \textbf{분산}이 큼
\item RF는 오른쪽(분산) 끝이 \textbf{거의 평탄}해서 ``많게 해두면
      된다'', GBDT는 \textbf{가파르게 올라서} ``정점을 찾아 멈춰야 한다''
\end{itemize}}
\end{frame}
